\begin{frame}{6.2 — Organização modular do monitor}
\begin{center}
\begin{tikzpicture}[node distance=0.65cm,every node/.style={draw,rounded corners,align=center,font=\small,text width=2.6cm}]
\node (main) {\texttt{main.c}\\cenário e impressão};
\node[right=of main,fill=green!10] (header) {\texttt{alarme.h}\\interface pública};
\node[right=of header] (source) {\texttt{alarme.c}\\implementação};
\draw[-{Latex}] (main)--node[above,draw=none,text width=1cm,font=\tiny]{inclui}(header);
\draw[-{Latex}] (source)--node[below,draw=none,text width=1cm,font=\tiny]{inclui}(header);
\end{tikzpicture}
\end{center}
\begin{itemize}
\item Cada arquivo tem uma responsabilidade identificável.
\item O cabeçalho informa como usar o módulo.
\item O arquivo \texttt{.c} contém a implementação da regra.
\item A interface permite trocar a implementação preservando o uso.
\end{itemize}
\end{frame}
